\chapter{Аналитическая часть}

В этом разделе будет дано определение электронной цифровой подписи, сформулированы алгоритмы создания и проверки цифровой подписи.




\section{Определение электронной цифровой подписи}

Электронная цифровая подпись (ЭЦП) --- это реквизит электронного документа, полученный в результате криптографического преобразования информации с использованием закрытого ключа подписи. ЭЦП позволяет подтвердить:
\begin{itemize}
    \item целостность документа (отсутствие изменений после подписания);
    \item принадлежность подписи владельцу (авторство);
    \item невозможность отказа от факта подписания (неотрекаемость).
\end{itemize}

Математически ЭЦП можно представить как результат работы функции подписи:
$$ S = \text{Sign}(M, K_{priv}) $$
где:
\begin{itemize}
    \item $M$ --- подписываемое сообщение (или его хеш);
    \item $K_{priv}$ --- закрытый ключ подписанта;
    \item $S$ --- полученная цифровая подпись.
\end{itemize}




\section{Алгоритмы создания цифровой подписи}

\subsection*{RSA-PSS (Probabilistic Signature Scheme)}
Схема вероятностной подписи на основе RSA:
\begin{enumerate}
    \item \textbf{Хеширование сообщения}: $h = Hash(M)$
    \item \textbf{Генерация случайной соли}: $salt$ (длиной обычно равной размеру хеша)
    \item \textbf{Формирование блока данных}: 
    $$ DB = Padding1 || salt $$
    \item \textbf{Вычисление маскированного блока}: 
    $$ maskedDB = DB \oplus MGF(H, len(DB)) $$
    где MGF --- функция генерации маски
    \item \textbf{Создание подписи}: 
    $$ S = (maskedDB || H || Padding2)^{d} \mod n $$
\end{enumerate}

\subsection*{ECDSA (Elliptic Curve Digital Signature Algorithm)}
Алгоритм на основе эллиптических кривых:
\begin{enumerate}
    \item \textbf{Параметры}: 
    \begin{itemize}
        \item Эллиптическая кривая $E$ над полем $F_p$
        \item Базовая точка $G$ порядка $n$
        \item Закрытый ключ $d$ ($1 < d < n-1$)
        \item Открытый ключ $Q = d \times G$
    \end{itemize}
    
    \item \textbf{Генерация подписи}:
    \begin{enumerate}
        \item Вычисление хеша: $h = Hash(M)$
        \item Выбор случайного $k$ ($1 < k < n-1$)
        \item Вычисление точки: $(x_1, y_1) = k \times G$
        \item $r = x_1 \mod n$ (если $r = 0$, повторить с другим $k$)
        \item $s = k^{-1}(h + r \cdot d) \mod n$ (если $s = 0$, повторить)
        \item Подпись: $(r, s)$
    \end{enumerate}
\end{enumerate}

\subsection*{EdDSA (Edwards-curve Digital Signature Algorithm)}
Современный алгоритм на основе кривых Эдвардса:
\begin{itemize}
    \item Использует кривые типа Curve25519 (Ed25519)
    \item Детерминированный --- не требует генерации случайных чисел
    \item Высокая производительность и безопасность
\end{itemize}




\section{Алгоритмы проверки цифровой подписи}

\subsection*{Проверка подписи RSA-PSS}
\begin{enumerate}
    \item \textbf{Восстановление подписи}: 
    $$ data = S^{e} \mod n $$
    \item \textbf{Разделение на компоненты}: $maskedDB || H || Padding2$
    \item \textbf{Восстановление блока DB}: 
    $$ DB = maskedDB \oplus MGF(H, len(maskedDB)) $$
    \item \textbf{Проверка формата}: убедиться, что DB имеет правильную структуру
    \item \textbf{Извлечение соли} и повторное вычисление хеша
    \item \textbf{Сравнение хешей}: если $H' = H$, подпись верна
\end{enumerate}

\subsection*{Проверка подписи ECDSA}
\begin{enumerate}
    \item \textbf{Проверка диапазона}: $r, s \in [1, n-1]$
    \item \textbf{Вычисление хеша}: $h = Hash(M)$
    \item \textbf{Вычисление промежуточных значений}:
    $$ w = s^{-1} \mod n $$
    $$ u_1 = h \cdot w \mod n $$
    $$ u_2 = r \cdot w \mod n $$
    \item \textbf{Вычисление точки}:
    $$ (x_1, y_1) = u_1 \times G + u_2 \times Q $$
    \item \textbf{Проверка}: 
    $$ r \equiv x_1 \mod n $$
\end{enumerate}

\subsection*{Общие принципы проверки ЭЦП}
\begin{itemize}
    \item \textbf{Аутентичность}: проверка соответствия открытому ключу
    \item \textbf{Целостность}: подтверждение неизменности документа
    \item \textbf{Своевременность}: проверка временных меток (при наличии)
    \item \textbf{Юридическая сила}: соответствие требованиям законодательства
\end{itemize}
